\section{Discussion}
The study supports removing one quarter of the GRABMyo monopolar recording sites within a defined mean recognition tolerance under this offline LOSO protocol. It does not show that every participant or gesture is preserved. The observed individual losses and corrected gesture declines make this distinction practically relevant: count selection for a real interface should include task priorities and user-specific failure criteria rather than relying solely on a pooled mean.

Attribution supplied predictive information beyond the mean random-subset baseline, and deletion with retraining supported its ranking relevance. Spatial coverage was a more competitive comparator. The matched primary experiment places geometric12's mean loss near the tolerance boundary; the direct mean advantage of SHAP12 remains uncertain. The full geometric curve required \RevGeomKStarCI\ retained sites under the original margin and decision rule, but only one deterministic geometric ordering was tested. These results motivate attribution as one evaluated selection procedure, with additional computational cost, rather than establishing a universal advantage over spatial selection. The shared-model experiment found a positive SHAP mean advantage over grouped permutation at six retained electrodes, but no established mean advantage at twelve. The gain under more aggressive reduction does not mean that six electrodes met the primary tolerance. It changes permutation grouping, participant-disjoint validation, scaler isolation and resources together, so it cannot attribute historical differences solely to computational budget.

Stability regularization was not an established improvement. Five-fold pooling obscures some person heterogeneity, which motivated the participant-level extension, but more descriptive dispersion does not imply that subtracting it improves classification. All penalty settings were fixed and disclosed rather than selected on outer scores. Correlated electrode features, normalization and shared attributing models complicate physiological interpretation; high attribution is model reliance and can reflect redundancy, task discrimination or montage geometry without identifying a particular muscle.

The inference has two distinct limits. First, the tolerance is an engineering choice, not a clinically validated equivalence bound; tighter or looser tolerances support different counts. Second, participant bootstrap intervals are conditional on fitted models with strongly overlapping training sets. The null-centred sensitivity corrects how the bootstrap test distribution is constructed, but cannot remove training-set dependence or historical development exposure. A nested development cohort, independently locked tolerance and external validation would strengthen future inference.

Several aspects restrict transfer. Only able-bodied public datasets were analyzed, not target clinical users. Zero-phase filtering is non-causal. Complete rest/window curves retain the primary ranking; the classifier sensitivity also retains it and samples only two counts. DB5 uses different sensors, sampling, labels and a small cohort, and derived bipolar pairs are constructed signals. The two split layouts differ and count selection had already seen the cohort. Full-array cross-session results do not demonstrate reduced-layout displacement robustness. No real-time accuracy, latency, energy, hardware cost reduction or clinical benefit was measured. Clinical channel-count experiments provide a further reason to keep this scope narrow: real-time prosthesis control can benefit from additional channels even when offline reduction appears promising \citep{simon2024implications}. Such evidence uses a different cohort, interface and endpoint and cannot validate or refute the present offline mean tolerance.

A prospective follow-up should fix a single geometry and count using development data only, record new participants with the documented montage, and define participant- and gesture-level loss criteria before testing. It should evaluate causal preprocessing and a closed-loop task, recording prediction latency and electrode failures as measured outcomes. Such a protocol is a proposal, not evidence supplied by the present datasets.

\section{Conclusion}
Cross-fitted attribution supported reduction of the GRABMyo monopolar array to \KStarMarginPtTwo\ electrodes at a 0.02 mean Macro-F1 tolerance under participant-independent offline evaluation. Utility beyond random selection was supported. Under shared resources, SHAP exceeded grouped permutation at six electrodes; its twelve-electrode mean advantage remained uncertain, as did the primary SHAP12--geometry12 contrast. Stability regularization did not establish an improvement. Participant losses, gesture-specific degradation, distinct fixed layouts and representation-dependent results restrict the claim to an evaluated reduction procedure. Independent validation is still required for a locked physical interface.
